\section*{Introduction}

Extracellular vesicles (EVs) are membrane-bound particles released through several routes, including multivesicular-body fusion, plasma-membrane budding and cell-disassembly pathways. Their abundance, molecular composition and biological activity depend on the state of the producer cell and on the mechanisms that generated them. Particle concentration alone is therefore an incomplete description of an EV preparation: populations that overlap in measured diameter can differ in origin, cargo and association with cellular injury \cite{welsh2024minimal}. This ambiguity is especially important when a physical stimulus is used to increase vesicle production, because the same perturbation may promote regulated secretion, repair-associated shedding or the release of damaged material.

Electric fields provide such a stimulus. Physiological-strength direct-current fields were recently shown to increase small-EV secretion by promoting intraluminal-vesicle formation and ceramide-dependent membrane budding \cite{huang2026electric}. At substantially shorter timescales, nanosecond pulsed electric fields (nsPEFs) create transient membrane permeability and redistribute ions without requiring sustained chemical stimulation \cite{beebe2003nanosecond,vernier2003calcium,semenov2013primary}. Measurements in urothelial cells further show that nsPEF susceptibility depends on cell size and that membrane resealing, osmotic swelling and stiffness loss depend on cell state and extracellular calcium \cite{kielbik2026selective}. These observations place membrane charging, calcium entry and membrane recovery upstream of the particle phenotype, but they do not establish how a nanosecond exposure becomes an extracellular response measured hours later.

Several biological processes could transmit this response across timescales. Calcium-dependent phospholipid scrambling rapidly changes the extracellularly exposed membrane proteome \cite{tuomivaara2024slapshot}, while TMEM16-family scramblases couple phosphatidylserine exposure to compensatory membrane internalization and repair \cite{kim2025tmem16e}. Extracellular ATP can stimulate EV secretion through P2X7-dependent calcium entry, linking mitochondrial state, purinergic signaling and vesicle release \cite{kao2023atp}. Membrane order and viscosity influence bending rigidity and therefore the energetic cost of vesicle formation \cite{steinkuhler2019mechanical}. Together, these studies support a causal sequence from electrical forcing to permeability, local calcium signaling, lipid and cytoskeletal remodeling, membrane repair and vesicle release. They also show why a model based on a single empirical relationship between field strength and particle yield would be difficult to transfer between pulse protocols or cell states.

A second difficulty lies in connecting intracellular mechanisms to particle measurements. Operational size classes do not uniquely identify biogenesis pathways, and extracellular concentration reflects both secretion and time-dependent losses caused by uptake, degradation, adsorption, aggregation, sampling and isolation. Measurements are further shaped by instrument response and recovery. A mechanistic release model therefore requires an explicit observation layer that maps overlapping pathway-derived size distributions to the measured window. Without this separation, observation-scale factors can be absorbed into biological rate constants and a decline in concentration can be misinterpreted as reduced secretion.

Here we developed and implemented a multiscale computational pipeline that represents pulse delivery, membrane and organelle electrodynamics, ionic and bioenergetic responses, remodeling and repair, pathway-resolved EV release, extracellular particle kinetics and measurement. The electrical pulse is compressed into physically interpretable descriptors, whereas downstream states are propagated over seconds to hours. We used the model to identify reversible, secretion-favored and injury-associated response regimes, and then connected simulated pathway output to size-resolved particle measurements from three pulse-labelled conditions and three harvest times. The fitted observation model reproduced the dominant measured size range but required an empirical exposure correction, revealing missing dose dependence rather than concealing it within pathway constants. Sensitivity and local-information analyses distinguished globally influential processes from parameters that were constrained by the current data and identified specific measurements needed for calibration.
